% =====================================================================
%  tr03b.tex —— 原文 p.9（tr-03）
%  内容：定义 7（Adams 运算）、定理 7（分裂原理）、§5 $K$ 理论的其他版本、
%        §5.1 实 $K$ 理论（定理 8）、§5.2 代数 $K$ 理论（定理 9、定义 8）
% =====================================================================
% ===== 原文 p.9 =====
% [原文 p.9]

该证明用到 $K$ 理论中 \textbf{Adams 运算}（Adams operations）的存在性以及 \textbf{分裂原理}（splitting principle）。

\begin{defn}{7}
  （Adams 运算）对紧空间 $X$ 与 $k\ge 0$，存在映射 $\psi^k:K(X)\to K(X)$，满足：
  \begin{itemize}[leftmargin=1.6em,itemsep=2pt]
    \item $f^*\psi^k=\psi^kf^*$，其中 $f:X\to Y$；
    \item 对线丛 $L$，$\psi^k(L)=L^k$；
    \item $\psi^k\psi^l=\psi^{kl}$；
    \item $\psi^p(x)\equiv x^p\pmod p$。
  \end{itemize}
\end{defn}
% [Tutor Note: 定义 7 的 (ii) $\psi^k(L)=L^k$ 与 (iii) $\psi^k\psi^l=\psi^{kl}$ 在线丛上自洽：由 (ii) 立刻得到 $\psi^k\psi^l(L)=L^{kl}=\psi^{kl}(L)$。对一般的 $E$，先用定理 7（分裂原理）把它沿 $f$ 拉回成线丛之和，再用 (i) 的函子性，$\psi^k$ 就被唯一确定——这正是需要分裂原理的原因。(iv) 里的 $p$ 指素数，它是 Frobenius 同余 $x^p\equiv x$ 在 $K$ 理论中的版本；注意右端 $x^p$ 是环 $K(X)$ 里的 $p$ 次幂，不是 $\psi$ 的复合。另：这里默认 $\psi^k$ 是环同态（从而保直和 $\psi^k(E\oplus F)=\psi^k(E)\oplus\psi^k(F)$），原文四条性质并未明列这一条，但下面的推理用到它。]

\begin{thm}{7}
  （分裂原理）给定紧空间上的向量丛 $E\xrightarrow{p}X$，存在紧空间 $P(E)$ 与映射 $P(E)\xrightarrow{f}X$，使得 $f^*P(E)$ 分裂为线丛之和，且 $f^*:K(X)\to K(F(E))$ 是单射。[这里 $P(E)$ 是 $X$ 上的射影丛（projective bundle）——从 $E\xrightarrow{p}X$ 中挖去零截面，再对 $\C^\times$ 作用取商；或者取每个纤维中的直线构成的空间。]
\end{thm}
% [Tutor Note: 定理 7 有两处原文笔误（1200 dpi 裁图逐字复核，另与文字层互证，印刷确为 "f^*P(E)" 与 "K(F(E))"）：(i)「$f^*P(E)$ 分裂为线丛之和」应为「$f^*E$ 分裂为线丛之和」——$P(E)$ 是空间（见方括号说明），能分裂成线丛之和的只能是底空间上的向量丛 $E$；(ii)「$K(F(E))$」应为「$K(P(E))$」——$f:P(E)\to X$ 的定义域（源空间）是 $P(E)$，故 $f^*$ 应看作 $K(X)\to K(P(E))$，而 $F(E)$ 在全文只出现这一次。按笔误译出是为了忠实；正确的陈述是：存在 $P(E)\xrightarrow{f}X$ 使 $f^*E\cong L_1\oplus\cdots\oplus L_n$ 且 $f^*:K(X)\to K(P(E))$ 为单射。]

\section{$K$ 理论的其他版本}

我们简要提及 $K$ 理论的其他化身（incarnations）的若干方面。

\subsection{实 $K$ 理论}

下面是实向量丛的 Bott 周期性版本。

\begin{thm}{8}
  存在同伦等价
  \[ \Omega^8BO\simeq BO\times\Z . \]
\end{thm}

因此与 $KO$ 相伴的上同调理论的周期是 8 而不是 2。
% [Tutor Note: 与定理 3 对照着读：复情形 $\Omega^2BU\simeq BU\times\Z$（周期 2），实情形 $\Omega^8BO\simeq BO\times\Z$（周期 8）。这里 $\Omega^8BO$ 是 $BO$ 的 8 重环路空间，由 $\Omega\Z\simeq*$ 可得 $\pi_{i+8}(BO)=\pi_i(BO)$（$i\ge 1$），这正是紧接上文那句「$KO$ 的周期是 8 而不是 2」的来源。]

\subsection{代数 $K$ 理论}

下述 Swan 定理显示了与其他数学领域的联系。

\begin{thm}{9}
  设 $X$ 是紧 Hausdorff 空间，$R=C(X;\R)$ 是 $X$ 上实值连续函数构成的环。若 $E\xrightarrow{p}X$ 是 $\R$-向量丛，则整体截面（global section）环 $\Gamma(X)=\{s:X\to E : p\circ s=1_X\}$ 是有限生成投射 $R$-模，且每个有限生成投射 $R$-模都以这种方式出现。
\end{thm}
% [Tutor Note: 定理 9 中 $1_X$ 是 $X$ 的恒等映射（$X\to X$），不是环 $R=C(X;\R)$ 的单位元；条件 $p\circ s=1_X$ 就是说 $s$ 保纤维，等价于 $s(x)\in E_x$。原文把截面空间记作 $\Gamma(X)$，通常写作 $\Gamma(X;E)$ 或 $\Gamma(E)$。最后一句「每个有限生成投射 $R$-模都以这种方式出现」是定理的要点：它把 $X$ 上的实向量丛与 $R$ 上的有限生成投射模一一对应，因此 $K_0(R)$ 正好是 $X$ 的实 $K$ 理论 $KO(X)$。]

聚焦于给定环上的有限生成投射模，就引向代数 $K$ 理论。

\begin{defn}{8}
  设 $R$ 是环，定义 $K_0(R)$ 为「有限生成投射 $R$-模在直和下构成的幺半群」的 Grothendieck 群。$K_0(R)$ 带有推广模张量积的乘法，因而成为一个环。$K_0$ 是从环范畴到环范畴的协变函子。每个环都有映射 $\Z\to R$，其诱导映射的余核给出约化群 $\widetilde{K}_0(R)$。
\end{defn}
% [Tutor Note: 定义 8 的约化群取的是 $\Z\to K_0(R)$ 的余核（cokernel），而拓扑情形里定义的是 $K(X)\to K(\ast)$ 的核（kernel），方向相反，不能照搬：代数里 $\Z\to K_0(R)$ 由 $1\mapsto[R]$（自由模 $R$）诱导，把它商掉才得到 $\widetilde{K}_0(R)$。若 $R$ 是数环，$\widetilde{K}_0(R)$ 就是 $R$ 的理想类群，这也正是原文接着说的 class group。]

例如，若 $R$ 是数环（number ring），则 $\widetilde{K}_0(R)$ 就是 $R$ 的类群（class group）。高阶 $K$ 群也有定义，但这里不再讨论。
